\section{Related Work}
\label{sec:related}

\paragraph{Structured pruning of LLMs.}
Width pruning methods (LLM-Pruner~\citep{llmpruner2023}, Bonsai~\citep{bonsai2025}, GRASPrune~\citep{graspprune2026}) remove attention heads or MLP channels while preserving dense computation.
Depth pruning~\citep{shortgpt2024,slicegpt2024,yang2024laco} removes entire layers but causes severe reasoning degradation~\citep{shrestha2026layerpruninglimits}; \citet{doingmore2026} revisit the trade-off under test-time scaling.
For reasoning models, RESP~\citep{resp2025}, RAC~\citep{rac2025}, and SSGR~\citep{ssgr2025} refine calibration data but retain gradient-based importance as the scoring criterion; \citet{llms2lrms2026} confirm width pruning is more robust than depth pruning for reasoning but do not study RLVR recovery.
A separate line uses functional structure to guide pruning: FANG~\citep{fang_2025} clusters FFN neurons by activation, FuncNet~\citep{funcnet_2025} decomposes LLMs via independent component analysis, and MPruner~\citep{mpruner2024} applies inter-layer CKA~\citep{kornblith2019cka} for depth pruning in CNNs.
We differ in both granularity and objective: \emph{intra-layer} CKA between individual structures, optimizing for RL recovery capacity rather than immediate accuracy.

\paragraph{Pruning-aware recovery.}
\citet{selfhealing2025} show that RLVR restores pruned reasoning models to near-original accuracy, outperforming SFT, but evaluate only importance-based unstructured pruning at low sparsity (96.5\% retention).
Under fine-grained structural pruning at 20--25\% sparsity, our controlled experiments indicate that this advantage diminishes when pruning preserves representational diversity, suggesting the effect is attributable to the pruning criterion rather than to RL itself.
Other work optimizes complementary axes: post-pruning scaling ($P^2$~Law;~\citealp{p2law2025}), recovery data selection (PASER;~\citealp{paser2025}), joint pruning-tuning (PAT;~\citealp{pat2025}), post-hoc attention compensation (RestoreLCC;~\citealp{restorelcc_2025}), iterative prune-tune loops~\citep{boiling_frog_2026}, post-RLVR compression (OPSD;~\citealp{opsd2026}), and rollout pruning during RLVR~\citep{xu2026arrol}---but none study which \emph{pruning criterion} enables RL-based recovery.
The Lottery Ticket Hypothesis~\citep{frankle2019lth} asks which subnetworks train effectively from scratch; our recovery capacity spectrum poses the analogous question for post-pruning recovery.
While LTH concerns pruning at initialization seeking sparse subnetworks that train from scratch, our setting involves post-training structured pruning followed by recovery; the analogy is conceptual (pruning strategy affects trainability) rather than technical.
Concurrent work benchmarks reasoning model compression~\citep{li2025reasoning_compression} and studies whether pruning can improve reasoning efficiency~\citep{canpruning2025}, but neither isolates the pruning criterion's effect on RL recovery.

\paragraph{RLVR training dynamics.}
\citet{rlsqueezes2026} show that RLVR \emph{compresses} incorrect reasoning trajectories while SFT \emph{expands} correct ones, predicting RL--SFT complementarity.
Our results qualify this prediction: the asymmetry manifests as a recovery capacity difference (the spectrum), but the expected accuracy gap is not observed when pruning preserves representational diversity ($|\Delta| \leq 1.8$\,pp across GSM8K, MBPP, and MATH-500; \S\ref{sec:decoupling}).
RL instability is well studied in the RL from Human Feedback (RLHF) context~\citep{gao2023overoptimization,dong2025rlplus,shah2025comedy}; \citet{mukherjee2025rlsubnetworks} find that RL fine-tuning updates only 5--30\% of parameters, raising the question of whether pruning criteria interact with this active subnetwork.
We connect these observations to structural pruning: the pruning criterion appears to shape susceptibility to RL collapse, while accuracy remains largely insensitive to the recovery paradigm at 20--25\% sparsity within our experimental scope.
